% Source: https://artofproblemsolving.com/community/c3h2077504_relatively_prime
% Stable identity: aops-topic-2077504-q01
Count the integers \(1,2,\dots,83\) divisible by \(2\), \(3\) or \(5\) by inclusion--exclusion:
\[
\left\lfloor\dfrac{83}{2}\right\rfloor+\left\lfloor\dfrac{83}{3}\right\rfloor+\left\lfloor\dfrac{83}{5}\right\rfloor-\left\lfloor\dfrac{83}{6}\right\rfloor-\left\lfloor\dfrac{83}{10}\right\rfloor-\left\lfloor\dfrac{83}{15}\right\rfloor+\left\lfloor\dfrac{83}{30}\right\rfloor=41+27+16-13-8-5+2=60.
\]
So \(83-60=\boxed{23}\) integers are relatively prime to \(30\), agreeing with the list.

The inclusion--exclusion attempt has an arithmetic slip: \(\left\lfloor\dfrac{84}{3}\right\rfloor=28\), not \(26\). Using \(84\) is harmless here because \(84\) itself is not coprime to \(30\): \(84-(42+28+16-14-8-5+2)=84-61=23\). The product formula \(N\left(1-\dfrac{1}{2}\right)\left(1-\dfrac{1}{3}\right)\left(1-\dfrac{1}{5}\right)\) is exact only when \(N\) is a multiple of \(30\) (it counts the residues coprime to \(30\) in complete blocks of \(30\)); \(84\) is not, so it only estimates the answer.
